% GENERATED FILE. Edit canonical lessons, then run npm run generate:books.
% canonical-source-hash: fnv1a64:e236ad95ec6a45f6
% canonical-lessons: ZH-C49-shangwu, ZH-W49-wu-noon, ZH-C49-zhongwu, ZH-C49-wanshang, ZH-W49-wan, ZH-C49-xiawu, ZH-C49-wanshanghao

\chapter{A Day from Morning to Night}
\label{ch:zh-a-day-from-morning-to-night}
\hlchaptermodality{49}

\begin{chapteropening}
\textbf{By the end of this chapter:} \emph{I can name the parts of the day and greet somebody in the evening.}

Complete the last lesson of chapter 49: I can name the parts of the day and greet somebody in the evening.
\end{chapteropening}

\section[\texorpdfstring{\zh{上午} (shàngwǔ)}{上午 (shàngwǔ)}]{\zh{上午} (shàngwǔ) — morning (before noon)}

\label{lesson:ZH-C49-shangwu}



\begin{quote}
Before the new one: say the Mandarin for how many, then the Mandarin for to wait.
\end{quote}

\subsection*{You'll want to know: \zh{上午}}
\textbf{\zh{上午}} — \emph{shàngwǔ} — \textquotedblleft{}morning (before noon)\textquotedblright{}.

\begin{quote}
\textbf{\zh{上午十点我去大学}} \emph{shàngwǔ shí diǎn wǒ qù dàxué} — At ten in the morning I go to the university.
\end{quote}

The new character, in pieces: \zh{午} — the top of \zh{年}: a short fall, two bars and a vertical. The next lesson puts it in your hand.

\begin{cousinweb}
\textbf{\zh{上午}} (\emph{shàngwǔ}) is built from \zh{上} 'upper, earlier' and \zh{午} 'noon', both inherited from Middle Chinese, \emph{dzyangH} and \emph{nguX} (Guangyun, in Baxter's transcription); the dictionary gives \emph{nguX} as the name of one of the branches of the old calendar.
\end{cousinweb}

\begin{grammarlens}[title={What you've built}]
One more word for times and partings.
\end{grammarlens}

\subsection*{Your turn}
\begin{itemize}
\raggedright
  \item \emph{Say it:} \emph{shàngwǔ}
  \item \emph{Say it:} \emph{shàngwǔ}, once more
  \item \emph{Say it:} \emph{shàngwǔ shí diǎn wǒ qù dàxué} — At ten in the morning I go to the university.
  \item \emph{Recall:} say the Mandarin for how many, then the Mandarin for to wait, then say \emph{shàngwǔ} again
\end{itemize}

\begin{tcolorbox}[breakable,colback=teal!4,colframe=teal!35!black,title={Before you move on}]
What does \zh{上午} mean? (\textquotedblleft{}Morning (before noon)\textquotedblright{}.) Say it once more.
\end{tcolorbox}
\section[\texorpdfstring{\zh{午} (wǔ)}{午 (wǔ)}]{\zh{午} — the top of \zh{年}}

\label{lesson:ZH-W49-wu-noon}



\begin{quote}
The word first, then the character inside it.

\begin{itemize}
\raggedright
  \item \emph{Recall:} say \emph{shàngwǔ} — \textbf{R1}, one lesson back, before it decays
\end{itemize}

\begin{quote}
\zh{年} \zh{十}
\end{quote}

\emph{Read it:} each one aloud, with nothing said first — \emph{nián}, \emph{shí}
\end{quote}

\subsection*{Script — the top of \zh{年}}
\hlblockfigure{\detokenize{figures/ZH-W49-wu-noon-filmstrip.pdf}}{How \zh{午} is written, stroke by stroke}

\begin{quote}
\zh{午}
\end{quote}

\textbf{One piece, four strokes.}

\textbf{A short left-falling stroke} at the top left.

\textbf{A shorter horizontal}, then \textbf{a long horizontal} under it.

\textbf{Last, the vertical} down through the middle of both bars.

It is \zh{年} with its short vertical and bottom bar left out.

\subsection*{Your turn}
\begin{itemize}
\raggedright
  \item \emph{Write it:} \zh{午}, all four, with the model in front of you
  \item \emph{Cover:} write it again from memory, and count the pieces, not the strokes
\end{itemize}

\begin{tcolorbox}[breakable,colback=teal!4,colframe=teal!35!black,title={Before you move on}]
Which character you know begins like \zh{午}? (\textbf{\zh{年}}.) How many strokes? (\textbf{Four}.) Which bar is longer? (\textbf{The lower one}.)

Source: \href{https://raw.githubusercontent.com/chanind/hanzi-writer-data/68d10a4b21150cae5e1ebbd223eed289cf32d90c/data/\%E5\%8D\%88.json}{Hanzi Writer Data \zh{午}, pinned stroke paths and medians}.
\end{tcolorbox}
\section[\texorpdfstring{\zh{中午} (zhōngwǔ)}{中午 (zhōngwǔ)}]{\zh{中午} (zhōngwǔ) — noon}

\label{lesson:ZH-C49-zhongwu}



\begin{quote}
Before the new one: say the Mandarin for to wait, then the Mandarin for before noon.
\end{quote}

\subsection*{You'll want to know: \zh{中午}}
\textbf{\zh{中午}} — \emph{zhōngwǔ} — \textquotedblleft{}noon\textquotedblright{}.

\begin{quote}
\textbf{\zh{中午我们吃饭}} \emph{zhōngwǔ wǒmen chī fàn} — We have lunch at noon.
\end{quote}

\begin{cousinweb}
\textbf{\zh{中午}} (\emph{zhōngwǔ}) is built from \zh{中} 'middle' and \zh{午} 'noon', both inherited from Middle Chinese, \emph{trjuwng} and \emph{nguX} (Guangyun, in Baxter's transcription).
\end{cousinweb}

\begin{grammarlens}[title={What you've built}]
One more word for times and partings.
\end{grammarlens}

\subsection*{Your turn}
\begin{itemize}
\raggedright
  \item \emph{Say it:} \emph{zhōngwǔ}
  \item \emph{Say it:} \emph{zhōngwǔ}, once more
  \item \emph{Say it:} \emph{zhōngwǔ wǒmen chī fàn} — We have lunch at noon.
  \item \emph{Recall:} say the Mandarin for to wait, then the Mandarin for before noon, then say \emph{zhōngwǔ} again
\end{itemize}

\begin{tcolorbox}[breakable,colback=teal!4,colframe=teal!35!black,title={Before you move on}]
What does \zh{中午} mean? (\textquotedblleft{}Noon\textquotedblright{}.) Say it once more.
\end{tcolorbox}
\section[\texorpdfstring{\zh{晚上} (wǎnshang)}{晚上 (wǎnshang)}]{\zh{晚上} (wǎnshang) — evening}

\label{lesson:ZH-C49-wanshang}



\begin{quote}
Before the new one: say the Mandarin for before noon, then the Mandarin for noon.
\end{quote}

\subsection*{You'll want to know: \zh{晚上}}
\textbf{\zh{晚上}} — \emph{wǎnshang} — \textquotedblleft{}evening\textquotedblright{}.

\begin{quote}
\textbf{\zh{晚上我看书}} \emph{wǎnshang wǒ kàn shū} — I read in the evenings.
\end{quote}

The new character, in pieces: \zh{晚} — \zh{日} on the left and a seven-stroke piece on the right. The next lesson puts it in your hand.

\begin{cousinweb}
\textbf{\zh{晚上}} (\emph{wǎnshang}) is built from \zh{晚} 'dusk' and the light ending shang, as in \zh{早上}. \zh{晚} is inherited from Middle Chinese \emph{mjonX} (Guangyun, in Baxter's transcription), glossed 'dusk'.
\end{cousinweb}

\begin{grammarlens}[title={What you've built}]
One more word for times and partings.
\end{grammarlens}

\subsection*{Your turn}
\begin{itemize}
\raggedright
  \item \emph{Say it:} \emph{wǎnshang}
  \item \emph{Say it:} \emph{wǎnshang}, once more
  \item \emph{Say it:} \emph{wǎnshang wǒ kàn shū} — I read in the evenings.
  \item \emph{Recall:} say the Mandarin for before noon, then the Mandarin for noon, then say \emph{wǎnshang} again
\end{itemize}

\begin{tcolorbox}[breakable,colback=teal!4,colframe=teal!35!black,title={Before you move on}]
What does \zh{晚上} mean? (\textquotedblleft{}Evening\textquotedblright{}.) Say it once more.
\end{tcolorbox}
\section[\texorpdfstring{\zh{晚} (wǎn)}{晚 (wǎn)}]{\zh{晚} — \zh{日}, then seven strokes}

\label{lesson:ZH-W49-wan}



\begin{quote}
The word first, then the character inside it.

\begin{itemize}
\raggedright
  \item \emph{Recall:} say \emph{wǎnshang} — \textbf{R1}, one lesson back, before it decays
\end{itemize}

\begin{quote}
\zh{日} \zh{多}
\end{quote}

\emph{Read it:} each one aloud, with nothing said first — \emph{rì}, \emph{duō}
\end{quote}

\subsection*{Script — \zh{日}, then seven strokes}
\hlblockfigure{\detokenize{figures/ZH-W49-wan-filmstrip.pdf}}{How \zh{晚} is written, stroke by stroke}

\begin{quote}
\zh{晚}
\end{quote}

\textbf{Two pieces, left before right.}

\textbf{\zh{日}, four strokes}, narrow on the left.

\textbf{The right piece, seven strokes.} A short left-falling stroke and a short stroke that runs across and falls, like the top of \zh{多}. A vertical, a horizontal that turns down, and a horizontal under it to close a small box. Then a left-falling stroke, and a vertical that bends right along the bottom and hooks up.

4 + 7 = eleven.

\subsection*{Your turn}
\begin{itemize}
\raggedright
  \item \emph{Write it:} \zh{晚}, all eleven, with the model in front of you
  \item \emph{Cover:} write it again from memory, and count the pieces, not the strokes
\end{itemize}

\begin{tcolorbox}[breakable,colback=teal!4,colframe=teal!35!black,title={Before you move on}]
How many pieces in \zh{晚}? (\textbf{Two} — \zh{日}, then seven strokes.) Which character you know starts the right piece the same way? (\textbf{\zh{多}}.) How does the last stroke end? (\textbf{It hooks up}.)

Source: \href{https://raw.githubusercontent.com/chanind/hanzi-writer-data/68d10a4b21150cae5e1ebbd223eed289cf32d90c/data/\%E6\%99\%9A.json}{Hanzi Writer Data \zh{晚}, pinned stroke paths and medians}.
\end{tcolorbox}
\section[xiàwǔ]{xiàwǔ — afternoon}

\label{lesson:ZH-C49-xiawu}



\begin{quote}
Before the new one: say the Mandarin for noon, then the Mandarin for evening.
\end{quote}

\subsection*{You'll want to know: xiàwǔ}
\emph{xiàwǔ} — \textquotedblleft{}afternoon\textquotedblright{}.

\begin{quote}
\emph{xiàwǔ wǒ hěn máng} — I'm busy in the afternoon.
\end{quote}

Its characters come later; for now it is a word for the ear.

\begin{cousinweb}
\emph{xiàwǔ} is built from xià 'lower, after' and \zh{午} 'noon', both inherited from Middle Chinese, \emph{haeX} and \emph{nguX} (Guangyun, in Baxter's transcription).
\end{cousinweb}

\begin{grammarlens}[title={What you've built}]
One more word for times and partings.
\end{grammarlens}

\subsection*{Your turn}
\begin{itemize}
\raggedright
  \item \emph{Say it:} \emph{xiàwǔ}
  \item \emph{Say it:} \emph{xiàwǔ}, once more
  \item \emph{Say it:} \emph{xiàwǔ wǒ hěn máng} — I'm busy in the afternoon.
  \item \emph{Recall:} say the Mandarin for noon, then the Mandarin for evening, then say \emph{xiàwǔ} again
\end{itemize}

\begin{tcolorbox}[breakable,colback=teal!4,colframe=teal!35!black,title={Before you move on}]
What does xiàwǔ mean? (\textquotedblleft{}Afternoon\textquotedblright{}.) Say it once more.
\end{tcolorbox}
\section[\texorpdfstring{\zh{晚上好} (wǎnshang hǎo)}{晚上好 (wǎnshang hǎo)}]{\zh{晚上好} (wǎnshang hǎo) — good evening}

\label{lesson:ZH-C49-wanshanghao}



\begin{quote}
Before the new one: say the Mandarin for evening, then the Mandarin for afternoon.
\end{quote}

\subsection*{You'll want to know: \zh{晚上好}}
\textbf{\zh{晚上好}} — \emph{wǎnshang hǎo} — \textquotedblleft{}good evening\textquotedblright{}.

\begin{quote}
\textbf{\zh{大家晚上好}} \emph{dàjiā wǎnshang hǎo} — Good evening, everyone.
\end{quote}

\begin{cousinweb}
\textbf{\zh{晚上好}} (\emph{wǎnshang hǎo}) is built from \zh{晚上} 'evening' and \zh{好} 'good', the same shape as \zh{早上好}.
\end{cousinweb}

\begin{grammarlens}[title={What you've built}]
One more word for times and partings.
\end{grammarlens}

\subsection*{Your turn}
\begin{itemize}
\raggedright
  \item \emph{Say it:} \emph{wǎnshang hǎo}
  \item \emph{Say it:} \emph{wǎnshang hǎo}, once more
  \item \emph{Say it:} \emph{dàjiā wǎnshang hǎo} — Good evening, everyone.
  \item \emph{Recall:} say the Mandarin for evening, then the Mandarin for afternoon, then say \emph{wǎnshang hǎo} again
\end{itemize}

\begin{tcolorbox}[breakable,colback=teal!4,colframe=teal!35!black,title={Before you move on}]
What does \zh{晚上好} mean? (\textquotedblleft{}Good evening\textquotedblright{}.) Say it once more.
\end{tcolorbox}
